\uuid{gEzo}
\titre{Injection/surjection}
\auteur{bodin ; hanani ; mzari}
\organisation{exo7}
\licence{CC BY-NC-SA 4.0}
\datecreate{2020-01-20}
\niveau{L1}
\module{Algèbre}
\chapitre{Logique, ensemble, raisonnement}
\sousChapitre{Ensemble}
\difficulte{1}
\theme{Ensembles, applications, injection, surjection}
\exo7qcmid{lille-1.ensembles.010}

\contenu{
\begin{question}
\qclassification{100.02,101.01,102.01,102.02}
\qtags{motcle=injection/surjection}
On considère l'application $f:\Nn\to \Nn$ définie par
$$\forall n\in \Nn,\; f(n)=n+1.$$
Quelle est la bonne réponse ?
\begin{answers}
    \bad{$f$ est surjective et non injective.}
    \good{$f$ est injective et non surjective.}
    \bad{$f$ est bijective.}
    \bad{$f$ n'est ni injective ni surjective.}
\end{answers}
\begin{explanations}
Si $f(n_1)=f(n_2)$ alors $n_1=n_2$, donc $f$ est injective. Par contre, $f(n)=0$ n'a pas de solution dans $\Nn$. Donc $f$ n'est pas surjective.
\end{explanations}
\end{question}
}
